\chapter{Conclusão}

Este trabalho apresentou o desenvolvimento de um protótipo funcional de sistema web para inventário patrimonial baseado em RFID, aplicado ao contexto do Colegiado de Ciência da Computação da UESC. A proposta partiu da identificação de limitações associadas ao controle patrimonial manual, como dependência de conferência individual, possibilidade de inconsistências e dificuldade de manter o inventário físico alinhado aos registros lógicos. Em relação ao pré-projeto, o núcleo da proposta foi preservado: automatizar parte do inventário patrimonial por meio de RFID, aplicação web e uma camada intermediária de integração entre dispositivos físicos e sistema.

O objetivo geral foi atendido por meio da implementação de uma aplicação composta por frontend, backend, API, banco de dados e camada de processamento de eventos RFID. O sistema permite cadastrar locais, leitores e itens patrimoniais, registrar leituras, atualizar localização física, executar auditorias, gerar logs e apontar inconsistências. Dessa forma, o protótipo demonstra uma base computacional viável para apoiar processos de inventário e rastreabilidade patrimonial.

Um aspecto central da solução é sua organização modular. A aplicação foi concebida para receber eventos de diferentes dispositivos, desde que estes se comuniquem com a API no formato esperado. Essa característica permite que a arquitetura evolua para cenários mais robustos, como o uso de sensores infravermelhos para acionar janelas de leitura e antenas RFID de maior alcance para captura automatizada de tags. Também permite integrar dispositivos simples, como leitores USB que operam como teclado ou periféricos semelhantes a leitores de código de barras, por meio de um comunicador intermediário. Na validação realizada, utilizou-se um leitor RFID de proximidade de baixo custo, escolha adequada à disponibilidade de equipamentos e ao objetivo de comprovar o fluxo lógico da solução.

Como limitação, destaca-se que a validação não avaliou desempenho de antenas RFID de longo alcance, leitura simultânea em massa, interferência eletromagnética, sensor infravermelho físico ou comportamento em ambiente institucional de produção. Também não se deve generalizar métricas de alcance ou acurácia sem ensaios específicos. O fluxo de acionamento por sensor foi preservado na arquitetura e verificado funcionalmente por software, mas sua validação física permanece como etapa futura. Assim, os resultados devem ser interpretados como validação funcional do software e do pipeline de processamento RFID, e não como validação completa de uma infraestrutura física definitiva.

Apesar dessas limitações, o protótipo contribui ao demonstrar que é possível integrar cadastro patrimonial, eventos RFID, auditoria, histórico operacional e tratamento de inconsistências em uma aplicação web. Como trabalhos futuros, recomenda-se testar a solução com antenas RFID de maior alcance, integrar sensores de presença como gatilhos de acionamento, documentar o hardware utilizado, ampliar a base de testes com maior quantidade de ativos, substituir o banco local por solução mais robusta, aprimorar autenticação e gerar relatórios formais de inventário e auditoria.

Conclui-se, portanto, que a solução desenvolvida atende ao propósito de validar uma arquitetura de software para inventário patrimonial baseado em RFID, oferecendo uma base extensível para futuras implantações e experimentos em ambientes acadêmicos. A principal evolução em relação à proposta inicial está na delimitação mais rigorosa da validação: o protótipo integrado foi entregue e testado com hardware acessível, enquanto as medições físicas de sensores, alcance, velocidade e leitura simultânea foram tratadas como continuidade experimental.
